\section*{M10: M9's upper envelope $U_3(\eta)$ is not attained}
\label{sec:k3-m10}

\textbf{Motivation.} M9 (\texttt{k3/general\_upper\_envelope.tex}) proves
$C_{3,\mathrm{ind}}^P(\eta)\le U_3(\eta)<2$ by bounding
$E_T^{\mathrm{proj}}\le T_1+T_2$ (triangle inequality) with
$T_1\le M^2q$, $T_2\le M\rho\sqrt{M^2-q^2}$ individually bounded via
operator-norm and closure-residual caps, then maximizing over
$q\in[0,\rho]$. This note asks the natural next question: can a single
configuration attain $U_3(\eta)$ exactly, i.e.\ can all of the
inequalities used in the derivation be simultaneously tight? The answer is
\textbf{no}, for a precise, provable, structural reason -- not merely
"not yet found."

\subsection*{A lemma about operator-norm-attaining directions}

\textbf{Lemma.} Let $N$ be linear (or a linear map restricted to a
2-dimensional subspace spanned by orthonormal $\hat u,\hat v$), with
$\|N(\cos\psi\,\hat u+\sin\psi\,\hat v)\|\le M$ for all $\psi$. If
$\|N(\hat u)\|=M$ exactly, then $N(\hat u)\perp N(\hat v)$.

\textit{Proof.} $f(\psi):=\|N(\cos\psi\,\hat u+\sin\psi\,\hat v)\|^2$
attains its maximum $M^2$ at $\psi=0$ (given). A smooth function on the
circle has zero derivative at an interior maximum, so
$f'(0)=2\langle N(\hat u),N(\hat v)\rangle\cos(0)\cdot(-\sin0)+\ldots=
2\langle N(\hat u),N(\hat v)\rangle=0$ (direct expansion of $f(\psi)=
\cos^2\psi\|N\hat u\|^2+\sin^2\psi\|N\hat v\|^2+2\sin\psi\cos\psi\langle
N\hat u,N\hat v\rangle$, differentiate, evaluate at $\psi=0$).
$\blacksquare$

This is a standard fact (a norm-attaining direction of a linear map is a
top right-singular vector, whose image is orthogonal to the image of any
right-singular vector for a different singular direction) -- the content
of this note is not the lemma itself but its two-fold application to the
k=3 chain's own telescoping structure.

\subsection*{Application at node 2}

Recall (\texttt{general\_upper\_envelope.tex}) $D_1\perp R_1$
(Pythagorean, at node 1), $\|D_1\|=q$, $\|R_1\|=\sqrt{M^2-q^2}$ (when
$\mu_1$ saturates its own bound $M$ at $(a,b)$ -- WLOG for an upper-bound
argument). Let $N(x):=\mu_2(x,c)$. Attaining $T_1=M^2q$ exactly requires
$\|S_1\|=\|N(D_1)\|=M\|D_1\|$, i.e.\ $N$ attains $M$ exactly at
$\hat u:=D_1/\|D_1\|$. By the Lemma (with $\hat v:=R_1/\|R_1\|$),
\[
  N(\hat u)\;\perp\;N(\hat v).
\]
Writing $p:=N(\hat u)$, $r:=N(\hat v)$: $S_1=qp$ and (choosing $P_2$ to
project onto $\mathrm{span}(p)$, which kills $p$ and preserves $r$ exactly
since $r\perp p$, without violating the closure-residual cap on
$(I-P_2)N$ -- see \texttt{k3\_non\_sharpness.py} for the explicit
construction) $S_2=\sqrt{M^2-q^2}\,r$. Since $p\perp r$:
\[
  \boxed{S_1\perp S_2 \quad\text{whenever $T_1$ is saturated exactly at node 2.}}
\]

\subsection*{Application at node 3 -- the contradiction}

Attaining $T_1=M^2q$ \emph{at the root} additionally requires
$\|\mu_3(S_1,d)\|=M\|S_1\|$, i.e.\ $L(x):=\mu_3(x,d)$ attains $M$ exactly
at $\hat s_1:=S_1/\|S_1\|$. Since $S_1\perp S_2$ (just shown), apply the
Lemma again with $\hat u=\hat s_1$, $\hat v=\hat s_2:=S_2/\|S_2\|$:
\[
  \mu_3(S_1,d)\;\perp\;\mu_3(S_2,d).
\]
But attaining the \emph{triangle-inequality} equality
$\|P_3\mu_3(S_1,d)+P_3\mu_3(S_2,d)\|=\|\mu_3(S_1,d)\|+\|\mu_3(S_2,d)\|$
(needed for $E_T^{\mathrm{proj}}=T_1+T_2$, the last step of the M9
derivation) requires $\mu_3(S_1,d)$ and $\mu_3(S_2,d)$ to be
\textbf{parallel, same-direction} (and both preserved by $P_3$).
Two nonzero orthogonal vectors are never parallel. \textbf{Contradiction.}

\subsection*{Conclusion}

\[
  \boxed{\text{$h(q)$ -- hence $U_3(\eta)$ -- is never exactly attained
  whenever both $T_1,T_2$ are nonzero at the optimum.}}
\]
The optimal $q^\star$ from M9 (either $q^\star=\rho$ or
$q^\star=M^2/\sqrt{M^2+\rho^2}$) is strictly positive for every
$\eta\in(0,1]$, and $T_2$'s coefficient $M\rho\sqrt{M^2-q^2}$ is nonzero
whenever $q<M$ (always true, since $q\le\rho\le M$ with equality only in
the degenerate $\eta=1,q=M$ boundary). So for every $\eta\in(0,1)$, both
terms are generically active at the optimum, and the obstruction applies:
\[
  \boxed{C_{3,\mathrm{ind}}^P(\eta) \;<\; U_3(\eta) \quad\text{strictly, for every } \eta\in(0,1).}
\]

\textbf{What this does not establish.} This is a negative result about
\emph{this specific decomposition-and-saturation mechanism} (the one M9's
own proof uses) -- it does not yet produce an exact tightened value
$U_{10}(\eta)<U_3(\eta)$. The natural next step is a genuine joint
optimization over the pair (magnitude allocation $q$, angle $\varphi$
between $S_1,S_2$), since \texttt{pareto\_frontier\_two\_direction\_norm\_
budget} (\texttt{k3\_non\_sharpness.py}) quantifies how the maximum
jointly-supportable magnitude pair shrinks as $\varphi\to0$ (full
alignment): at $\varphi=\pi/2$ the joint cap recovers the independent
per-term caps used by M9 (i.e.\ $U_3$'s own derivation), while at
$\varphi=0$ the cap drops to $M/\sqrt{1+\text{ratio}^2}$ -- a strictly
smaller number whenever the ratio is nonzero. \textbf{This joint
$(q,\varphi)$ optimization was set up but not solved this session}; it is
recorded as the precise open problem, not vaguely as "improve the bound."

\textbf{A numerical search was also attempted and was inconclusive.} A
gradient-based (Adam + penalty method) adversarial search over raw
tensors (\texttt{m10\_k3\_search.py}, run this session, not committed to
the repository because its results are not trustworthy) failed to even
recover the already-known certified lower witness $L_3(\eta)$ at any
tested point -- e.g.\ at $\eta=0.5$ it found a normalized value of
$0.51$ against a known-achievable $L_3(0.5)=1.73$. This indicates a
methodology problem (poor differentiable operator-norm surrogates via
power iteration, bad initialization scale, or insufficient steps), not a
mathematical finding, and is reported here only so a future session does
not repeat the same broken setup without knowing it was already tried and
discarded.

\textbf{Status.} \texttt{PROVED} (the non-sharpness of $U_3(\eta)$, via
the two-fold Lemma application above). The exact tightened envelope
remains \texttt{OPEN} -- a well-posed joint optimization, not an
unscoped question.
